\documentclass{ximera}


\title{Isaiah Hayes - 21 April 2026 - Project 3: Kaleidoscope Project}

\begin{document}


\begin{center}
    This project explores the mathematical construction of a virtual kaleidoscope using linear transformations, matrix representations, and composition.
\end{center}

A physical kaleidoscope starts with a pie-shaped image. Two mirrors are held perpendicular to the paper with the image: one mirror is held above the horizontal edge, and the other is held above the angled edge. These mirrors meet at a $60^\circ$ angle.

We represent the coordinates of each point in the original image using two-dimensional vectors $\vec{x} = \langle x_1, x_2 \rangle$. After a reflection, the new coordinates $\vec{y} = \langle y_1, y_2 \rangle$ are determined by a linear transformation:
\[
\begin{aligned}
    y_1 &= a_{11}x_1 + a_{12}x_2 \\
    y_2 &= a_{21}x_1 + a_{22}x_2
\end{aligned}
\]
This can be expressed in matrix form as:
\[
\begin{bmatrix} y_1 \\ y_2 \end{bmatrix} = \begin{bmatrix} a_{11} & a_{12} \\ a_{21} & a_{22} \end{bmatrix} \begin{bmatrix} x_1 \\ x_2 \end{bmatrix} \quad \text{or} \quad \vec{y} = A\vec{x}
\]

\textbf{Why are there only six pie-shaped pieces or sectors?}

There are exactly six sectors because the mirrors are set at a $60^\circ$ angle. Since a complete revolution around the origin is $360^\circ$, the mirrors divide the plane into:
\[
\frac{360^\circ}{60^\circ} = 6 \text{ equal sectors.}
\]
Mathematically, these sectors are generated by repeatedly applying the reflection and rotation transformations. Because 60 is a divisor of 360, the reflections eventually "close" the circle perfectly without any overlapping or gaps, creating a symmetric hexagonal pattern.

\textbf{(1) Reflection across the x-axis}

To reflect an image across the x-axis, the horizontal position ($x_1$) remains the same while the vertical position ($x_2$) is negated. This transformation maps the basis vector $\mathbf{e}_1 = \langle 1,0 \rangle$ to $\langle 1,0 \rangle$ and $\mathbf{e}_2 = \langle 0,1 \rangle$ to $\langle 0,-1 \rangle$. The resulting matrix is:
\[
\begin{bmatrix}
\answer{1} & \answer{0} \\
\answer{0} & \answer{-1}
\end{bmatrix}
\]

\textbf{(2) Rotation by 60 Degrees}

A counterclockwise rotation by an angle $\theta$ is represented by the standard rotation matrix $R = \begin{bmatrix} \cos\theta & -\sin\theta \\ \sin\theta & \cos\theta \end{bmatrix}$. For a $60^\circ$ rotation, we use $\cos(60^\circ) = 1/2$ and $\sin(60^\circ) = \sqrt{3}/2$. The matrix is:
\[
\begin{bmatrix}
\answer{1/2} & \answer{-\sqrt{3}/2} \\
\answer{\sqrt{3}/2} & \answer{1/2}
\end{bmatrix}
\]

\textbf{(3) Reflection across the 60-degree diagonal}

To find the matrix for a reflection across a line at an angle $\theta = 60^\circ$, we use the general diagonal reflection matrix formula:
\[
\begin{bmatrix} \cos(2\theta) & \sin(2\theta) \\ \sin(2\theta) & -\cos(2\theta) \end{bmatrix}
\]
Since the mirror is at $60^\circ$, we calculate the values for $2(60^\circ) = 120^\circ$. We know that $\cos(120^\circ) = -1/2$ and $\sin(120^\circ) = \sqrt{3}/2$. Therefore, the matrix is:
\[
\begin{bmatrix}
\answer{-1/2} & \answer{\sqrt{3}/2} \\
\answer{\sqrt{3}/2} & \answer{1/2}
\end{bmatrix}
\]

\subsection*{Derivation of the Reflection Matrix}
I developed this matrix by applying the double-angle reflection formula. This transformation flips a point across the line $y = \tan(60^\circ)x$. By considering the transformation of the basis vectors, we see that the line of symmetry effectively swaps and scales the coordinates based on the $120^\circ$ orientation.

\end{document}